\section{Twisted intersection forms and period matrices}\label{sec:periods}
Our goal in this section and the next is to compute the determinant of 
$(V_b^{-1})^{\cB}$ in the realizable case.  In Section~\ref{subsec:betti-pairing} we recall Lam's result expressing the inverse of the Varchenko matrix in terms of his combinatorial Betti homology matrix.  Section~\ref{subsec:deRham-pairing} recalls Lam's combinatorial de Rham cohomology pairing and determinant formula for the pairing with respect to the weighted $\beta$nbc basis.  Section~\ref{subsec:geometric-comparison} summarizes Lam's identification of the combinatorial pairings with geometric pairings.  Section~\ref{subsec:period} introduces period matrices and Lam's version of the twisted period relation of Cho--Matsumoto.  In Section~\ref{subsec:period-determinant} we recall Varchenko, Loeser--Sabbah, Douai--Terao's determinant formula for the period matrices.  All of these ingredients are finally combined in Section~\ref{sec:determinants} to complete our proof of Theorem~\ref{thm:Lam-real}.


\subsection{Lam's combinatorial Betti pairing}\label{subsec:betti-pairing}
Let $(\M,0)$ be an affine oriented matroid of rank $r+1$ with ground set $E$, and let
$b_i$, $i\in E$, be independent variables, including $b_0$.
Lam's \emph{Betti homology matrix}
\[
 B_b=\bigl(B_b( P,Q) \bigr)_{P,Q\in\cF}
\]
is a symmetric matrix over $\QM(b_i\mid i\in E)$, defined by a sum over
flags of faces \cite[Definition 8.2]{Lam}. Its summands are signed
products of monomials $b_F$ and factors $(b_F^2-1)^{-1}$, where $F$
ranges over zero sets of nonzero faces of $P$. We will only need this
feature of the definition and the following formula for the inverse.

\begin{proposition}[{\cite[Theorem 9.2 and Corollary 9.3]{Lam}}]\label{prop:Betti-Varchenko}
For $P,Q\in\cF$,
\[
 (B_b^{-1})(P,Q)=\frac{(-1)^r}{1-b_E}
 \bigl(b_{\Sep(P,Q)}+(-1)^{r+1}b_{E\setminus\Sep(P,Q)}\bigr).
\]
In particular,
\[
 \left.B_b\right|_{b_0=0}=(-1)^rV_b^{-1}.
\]
\end{proposition}


We are interested in the $\cB \times \cB$-submatrix $B_b^{\cB}$.
This block is independent of $b_0$. Indeed, every nonzero face of a
bounded tope is feasible, so none of the zero sets in its face flags
contains $0$. Proposition~\ref{prop:Betti-Varchenko} therefore gives
\[
 B_b^{\cB}=\left( \left.B_b\right|_{b_0=0}\right)^\cB = (-1)^r(V_b^{-1})^{\cB}.
\]
Since $|\cB|=\beta(M)$, the determinants differ by the sign
\[
 \det B_b^{\cB}=(-1)^{r\beta(M)}
                 \det\bigl((V_b^{-1})^{\cB}\bigr).
\]
\begin{remark}
Theorem~\ref{thm:Lam-real} is stated in this Varchenko normalization.
The Betti determinant in \cite[Conjecture 8.10]{Lam} is missing this global sign, which is already visible
in \cite[Example 8.11]{Lam}.
\end{remark}
\begin{example}\label{ex:line}
For the two-point arrangement defined by $f_1(x)=x$ and $f_2(x)=1-x$,
the unique bounded chamber is $P=(0,1)$, oriented from $0$ to $1$.
All matrices in this example have size $1\times1$.
The Betti self-intersection of $P$ is
\[
 \det B_b^{\cB}=B_b(P,P)=1+\frac{1}{b_1^2-1}+\frac{1}{b_2^2-1}
     =-(V_b^{-1})(P,P) = -\det\bigl((V_b^{-1})^{\cB}\bigr),
\]
as in \cite[Example 14.2]{Lam}. Here the sign in
Proposition~\ref{prop:Betti-Varchenko} is $(-1)^r=-1$.
Example~\ref{ex:line4} recovers this entry from the
period relation.
\end{example}
\subsection{Lam's combinatorial de Rham pairing}\label{subsec:deRham-pairing}
Let $(M,0)$ be an affine matroid of rank $r+1$ with ground set $E = E_0 \sqcup \{0\}$, and let $a_i$, $i\in E$ be independent variables.
Let $A^{\bullet}(M)$ be the Orlik--Solomon algebra over $\CM$.
It is the exterior algebra on generators $e_i$, $i\in E$, modulo
the elements $\partial e_C$ for circuits $C$, where $\partial$ is
the derivation of degree $-1$ satisfying $\partial e_i=1$, and has top
degree $r+1$.
The reduced algebra is $\bar A^{\bullet}(M)=\ker\partial$.
It is generated by $\bar e_i=e_i-e_0$, $i\in E_0$, has top
degree $r$ and $A^{r+1}(M) \cong \bar A^r(M)$.  

Lam's combinatorial de Rham cohomology pairing is a bilinear form on
$\bar A^r(M)$, the top part of the reduced Orlik--Solomon algebra,
with values in rational functions in the variables $a_i$, $i\in E$:
\[
 \langle~,~\rangle^{\dR}:
 \bar A^r(M)\otimes_{\CM}\bar A^r(M)
 \longrightarrow\CM(a_i\mid i\in E).
\]
The bilinear form is defined by residues. For a flag of flats
\[
 F_\bullet=(\varnothing=F_0\subsetneq F_1\subsetneq\cdots
             \subsetneq F_r\subsetneq E),\qquad \rk F_j=j,
\]
let $\Res_{F_\bullet}:\bar A^r(M)\to\CM$ be the iterated
Orlik--Solomon residue of \cite[Definition 5.2]{Lam}. It takes residues
successively for the rank-one flats $F_j/F_{j-1}$ in $M/F_{j-1}$.
For $x,y\in\bar A^r(M)$, Lam's bilinear form is
\begin{equation}\label{eq:combinatorial-deRham}
 \langle x,y\rangle^{\dR}
 =\sum_{F_\bullet}
   \frac{\Res_{F_\bullet}(x)\Res_{F_\bullet}(y)}
        {a_{F_1}\cdots a_{F_r}},
\end{equation}
where the sum runs over all complete flags and $a_F = \sum_{i\in F} a_i$. This is \cite[Definition 5.6 and Corollary 5.12]{Lam} on the reduced
algebra.

This bilinear form descends to give a bilinear form on Aomoto cohomology.  To define Aomoto cohomology, we specialize the variables $a_i, i \in E$ to complex numbers such that $\sum_{i\in E} a_i = 0$.  The Aomoto cohomology $A_{\mathbf a}$ depends on the choice $\mathbf a=(a_i)_{i\in E}$ of specialization.  For any $J \subseteq E$, let $a_J=\sum_{i\in J}a_i$.  In addition to the assumption that $a_E=\sum_{i\in E} a_i =0$ (equivalently, $a_0=-\sum_{i\in E_0}a_i$) we assume that the weights $a_i, i\in E$ are \emph{nonresonant}, meaning
\begin{equation}\label{eq:weights}
 a_F\notin\ZM\quad
 \text{ for all }F\in\cL\setminus\{\varnothing,E\}\text{ with }\ M^F\text{ connected}.
\end{equation}
While the denominators in the expression for $\langle x,y\rangle^{\dR}$ above are products ${a_{F_1}\cdots a_{F_r}}$ for arbitrary flags, Lam~\cite[Corollary 10.25]{Lam} shows that $\langle x,y\rangle^{\dR}$ can be re-expressed such that only products of $a_F$ for \emph{connected} flats $F \neq \hat 0, \hat 1$ appear in the denominator.  In particular, under the condition (\ref{eq:weights}), the rational function $\langle x,y\rangle^{\dR}$ on $a_E=0$ is regular.
The nonresonance condition is also used in the basis statements below and the geometric comparisons in
Section~\ref{subsec:geometric-comparison}.

Put $\omega_{\mathbf a}=\sum_{i\in E_0}a_i\bar e_i$.
Equipping $\bar A^{\bullet}(M)$ with differential
$\omega_{\mathbf a}\wedge(-)$ gives the \emph{Aomoto complex}.
The \emph{Aomoto cohomology} is defined as
\[
 A_{\mathbf a}
 :=H^r(\bar A^{\bullet}(M),\omega_{\mathbf a}\wedge(-))
 =\frac{\bar A^r(M)}{\omega_{\mathbf a}\wedge\bar A^{r-1}(M)}.
\]
Thus $A_{\mathbf a}=A_{-\mathbf a}$ as quotients of $\bar A^r(M)$.
We retain the two labels to distinguish their later geometric
realizations with dual local systems.
Under \eqref{eq:weights}, the Aomoto cohomology is concentrated in
degree $r$ and has dimension $\beta(M)$ \cite[Theorem 7.3]{Lam}.

At nonresonant weights, Lam's de Rham cohomology pairing $\langle~,~\rangle^{\dR}$ descends to a nondegenerate symmetric pairing on $A_{\mathbf a}$, which we denote by $\langle~,~\rangle^{\dR}_{\mathbf a}$ 
\cite[Theorem 7.7 and Corollary 7.15]{Lam}.

We now recall the \textit{weighted $\beta$nbc basis} of $A_{\mathbf a}$. Fix a total order on $E$ with $0$ least.
A broken circuit is a circuit
with its least element removed. An nbc basis contains no broken
circuit. It is a $\beta$nbc basis if for every $i\in B\setminus\{0\}$
there is $j<i$ such that $(B\setminus\{i\})\cup\{j\}$ is again a
basis. Write $\Bnbc$ for this set. Its cardinality is $\beta(M)$
\cite{Ziegler1992}, \cite[Definition 7.9 and Theorem 7.10]{Lam}.

Write a $\beta$nbc basis in decreasing order as
$B=(i_1>\cdots>i_r>i_{r+1}=0)$ and put
\[
 F_j(B)=\cl_M\{i_1,\ldots,i_j\},\qquad
 \omega_{\mathbf a}(F)=\sum_{i\in F}a_i\bar e_i.
\]
Here $0\notin F_j(B)$. The corresponding weighted form is
\begin{equation}\label{eq:nbc-form}
 \psi_{\mathbf a}^B
 =\omega_{\mathbf a}(F_r(B))\wedge\cdots\wedge
   \omega_{\mathbf a}(F_1(B)).
\end{equation}
At nonresonant weights these forms give a basis of $A_{\mathbf a}$
\cite{FalkTerao1997}, \cite[Corollary 7.14]{Lam}.
They satisfy $\psi_{-\mathbf a}^B=(-1)^r\psi_{\mathbf a}^B$.

Set
\[
 D_{\mathbf a}^{\mathrm{nbc}}(B,B')
   =\langle\psi_{\mathbf a}^B,\psi_{-\mathbf a}^{B'}\rangle^{\dR}_{\mathbf a}.
\]
This uses the basis at its own weight in each argument. Relative
to Lam's Gram matrix using $\psi_{\mathbf a}$ in both arguments,
$D_{\mathbf a}^{\mathrm{nbc}}$ is multiplied by $(-1)^r$.

Recall $w_F=\beta(M^F)\beta(M_F)$ and put
\[
 R_M(\mathbf a)
 =\frac{\displaystyle\prod_{F\in\cL_0\setminus\{E\}}a_F^{w_F}}
        {\displaystyle\prod_{F\in\cL\setminus(\cL_0\cup\{\varnothing\})}
             a_F^{w_F}}.
\]

\begin{theorem}[Lam]\label{thm:deRham-determinant}
As rational functions on $a_E=0$,
\[
 \det D_{\mathbf a}^{\mathrm{nbc}}=\pm R_M(\mathbf a)^{-1}.
\]
The sign is constant.
\end{theorem}
\begin{proof}
This is \cite[Corollary 7.20]{Lam}.
The use of $\psi_{-\mathbf a}$ in the second argument changes the
determinant by the constant sign $(-1)^{r\beta(M)}$.
\end{proof}

\begin{example}\label{ex:line2}
Consider again the two-point arrangement from Example~\ref{ex:line}.
Its matroid $M$ is the uniform matroid $U_{2,3}$ on $E=\{0,1,2\}$.
For the order $0<1<2$, the nbc bases are $\{0,1\}$ and $\{0,2\}$,
and the unique $\beta$nbc basis is $B=\{0,2\}$.
The weighted form is $\psi^B_{\mathbf a}=a_2\bar e_2$.
Since $\psi^B_{-\mathbf a}=-\psi^B_{\mathbf a}$,
\eqref{eq:combinatorial-deRham} and $a_0=-a_1-a_2$ give
\[
 \det D_{\mathbf a}^{\mathrm{nbc}} = D_{\mathbf a}^{\mathrm{nbc}}(B,B)
 =-a_2^2\left(\frac1{a_2}+\frac1{a_0}\right)
 =\frac{a_1a_2}{a_0}.
\]
Note that this agrees with Lam's Theorem~\ref{thm:deRham-determinant}.
\end{example}

\begin{remark}\label{rem:chamber-deRham}
For an orientation $\M$ of $M$, Eur and Lam associate to each tope $P$ a weight-independent canonical
element $\Omega_P\in A^{r+1}(M)$ \cite[Theorem 2.11]{EurLam}.
Put $\bar\Omega_P=\partial\Omega_P$.
At nonresonant weights their classes for $P\in\cB$ form a second
basis of $A_{\mathbf a}$ \cite[Corollary 7.14]{Lam}. Its Gram matrix
\[
 D_{\mathbf a}^{\mathrm{ch}}(P,Q)
 =\langle\bar\Omega_P,\bar\Omega_Q\rangle^{\dR}_{\mathbf a}
\]
has determinant $\pm R_M(\mathbf a)$ by
\cite[Theorem 7.18]{Lam}.

However, we focus our attention on the Gram matrix for the weighted $\beta$nbc basis (instead of the canonical chamber basis) because, as recalled in Section~\ref{subsec:period-determinant}, the determinants of the corresponding period matrices are known.
\end{remark}

\subsection{Geometric realizations and comparison}\label{subsec:geometric-comparison}
We will now recall Lam's geometric realizations of the Betti and de Rham matrices defined above. For this, in this section we assume that $\C$ is a real essential affine arrangement
and the weights $a_i$ are nonresonant, that is, \eqref{eq:weights} holds.
Let
\[
 U=\CM^r\setminus\bigcup_{i\in E_0}H_i^{\CM}
\]
be the complement of the complexified arrangement. Specialize
$b_i=e^{-\pi i a_i}$, so $b_E=1$. We use the bounded block
$B_b^{\cB}$, which is independent of $b_0$.

The logarithmic connection
\[
 \nabla_{\mathbf a}=\dd+\omega_{\mathbf a}\wedge(-),\qquad
 \omega_{\mathbf a}=\sum_{i\in E_0}a_i\,\dd\log f_i
\]
has a rank-one local system $\Loc_{\mathbf a}$ of flat sections.
Locally these are scalar multiples of
$\varphi_{\mathbf a}^{-1}=\prod_{i\in E_0}f_i^{-a_i}$.
The monodromy around $H_i$ is $b_i^2$ and
$\Loc_{\mathbf a}^{\vee}=\Loc_{-\mathbf a}$.

We use ordinary and Borel--Moore homology, and ordinary and compactly
supported cohomology, with coefficients in these local systems.
Under \eqref{eq:weights}, the natural maps
\[
 H_r(U,\Loc_{\pm\mathbf a})\longrightarrow
 H_r^{\BM}(U,\Loc_{\pm\mathbf a}),\qquad
 H_c^r(U,\Loc_{\pm\mathbf a})\longrightarrow
 H^r(U,\Loc_{\pm\mathbf a})
\]
are isomorphisms by \cite[Theorem 12.3]{Lam}. Kohno proved this under
stronger nonresonance assumptions in \cite[Theorem 1]{Kohno}.
We denote their inverses by $\reg$.
The natural maps, and hence their inverses, are dual under the perfect
evaluation pairings \cite[equation (12.3)]{Lam}. Explicitly,
\[
 (\xi,\reg c)=(\reg\xi,c),
 \qquad \xi\in H^r(U,\Loc_{-\mathbf a}),\quad
 c\in H_r^{\BM}(U,\Loc_{\mathbf a}).
\]
The left side pairs cohomology with ordinary homology, and the right
side pairs compactly supported cohomology with Borel--Moore homology.

These groups vanish outside degree $r$, where they have dimension
$\beta(M)$ \cite[Theorem 12.3 and Section 12.3]{Lam}.
The algebraic de Rham groups are
\[
 H^\bullet_{\dR}(U,\nabla_{\pm\mathbf a})
 =H^\bullet(\Gamma(U,\Omega_U^\bullet),\nabla_{\pm\mathbf a}).
\]
The de Rham comparison gives isomorphisms
\[
 \comp_{\pm\mathbf a}:H^\bullet_{\dR}(U,\nabla_{\pm\mathbf a})
 \xrightarrow{\ \sim\ }H^\bullet(U,\Loc_{\pm\mathbf a})
\]
\cite[Section 12.3]{Lam}. We also write 
\[\reg: H^r_{\dR}(U,\nabla_{\pm\mathbf a}) \to H^r_{c,\dR}(U,\nabla_{\pm\mathbf a}) \] for the corresponding
regularization in de Rham cohomology.

Consider the logarithmic realization
\[
 \iota:\bar A^{\bullet}(M)\hookrightarrow\Gamma(U,\Omega_U^{\bullet}),
 \qquad \bar e_i\longmapsto\dd\log f_i=\frac{\dd f_i}{f_i},
\]
which induces a map of complexes
\[
 (\bar A^{\bullet}(M),\omega_{\mathbf a}\wedge(-))
 \xhookrightarrow{\ \iota\ }
 (\Gamma(U,\Omega_U^{\bullet}),\nabla_{\mathbf a}).
\]
Under \eqref{eq:weights}, this is a quasi-isomorphism by the theorem
of Esnault--Schechtman--Viehweg \cite{EsnaultSchechtmanViehweg1992}
and Schechtman--Terao--Varchenko \cite{SchechtmanTeraoVarchenko1995},
as stated in \cite[Theorem 12.4]{Lam}.

Taking top cohomology gives the identifications
\[
 A_{\pm\mathbf a}
 \xrightarrow{\ \iota_*\ }H^r_{\dR}(U,\nabla_{\pm\mathbf a})
 \xrightarrow{\ \comp_{\pm\mathbf a}\ }H^r(U,\Loc_{\pm\mathbf a}).
\]
The first sends $[\eta]$ to the twisted de Rham class $[\iota(\eta)]$.
The second is the de Rham comparison with the local system of flat
sections.
Write $\per_{\pm\mathbf a}=\comp_{\pm\mathbf a}\circ\iota_*$
for the composite.
Under $\iota$, the elements $\bar\Omega_P$ become the canonical
rational top forms of the chamber closures, with logarithmic poles
along their boundaries \cite[Remark 4.10]{Lam}.

Choose once and for all the branch $\log$ with cut along the positive
imaginary ray, agreeing with the real logarithm on $\R_{>0}$. We write
\[
 z^a=\exp(a\log z).
\]
The product
\[
 \varphi_{\mathbf a}
 =\prod_{i\in E_0}f_i^{a_i}
\]
is holomorphic on the complex neighbourhood $P+i\R^r$ of every real
chamber $P$.
Put $P^-=\{i\in E_0\mid P(i)=-\}$ and define the normalized branches by
\begin{equation}\label{eq:loading}
 \varphi_{\mathbf a,P}
 =b_{P^-}^{-1}\varphi_{\mathbf a},\qquad
 \varphi_{\mathbf a,P}^{-1}=b_{P^-}\varphi_{\mathbf a}^{-1}.
\end{equation}
Its restriction to the real chamber satisfies
\[
 \varphi_{\mathbf a,P}(x)=\prod_{i\in E_0}|f_i(x)|^{a_i},
 \qquad x\in P.
\]
Thus it is positive on $P$ for real weights $\mathbf a$, and
$\varphi_{-\mathbf a,P}=\varphi_{\mathbf a,P}^{-1}$.
The chosen branch extends by analytic continuation along paths in $U$.
This is the normalization of loaded chambers in
\cite[Definition 12.6 and equation (12.5)]{Lam}.

With the ambient orientation of $\R^r$, the loaded chambers define
classes in Borel--Moore homology
\begin{align*}
 C_{\mathbf a}(P)&=[P\otimes\varphi_{\mathbf a,P}^{-1}]
       \in H_r^{\BM}(U,\Loc_{\mathbf a}),\\
 C_{-\mathbf a}(P)&=[P\otimes\varphi_{\mathbf a,P}]
       \in H_r^{\BM}(U,\Loc_{-\mathbf a}).
\end{align*}
The classes $C_{\mathbf a}(P)$ and $C_{-\mathbf a}(P)$, for $P\in\cB$, are bases
of the respective Borel--Moore groups
\cite{Kohno}, \cite[Proposition 12.7]{Lam}.

The Betti homology pairing
\[
 \langle-,-\rangle_{\mathrm{hom}}:
 H_r^{\BM}(U,\Loc_{-\mathbf a})\otimes
 H_r^{\BM}(U,\Loc_{\mathbf a})\longrightarrow\CM
\]
is obtained by applying $\reg$ to the second argument and then the
Poincar\'e--Verdier duality pairing \cite[Definition 12.5]{Lam}.
Using regularization to identify Borel--Moore and ordinary homology,
its dual under the perfect evaluation pairings is the Betti
cohomology pairing
\[
 \langle-,-\rangle_{\mathrm{coh}}:
 H^r(U,\Loc_{\mathbf a})\otimes H^r(U,\Loc_{-\mathbf a})
 \longrightarrow\CM
\]
\cite[equation (12.3) and Section 12.4]{Lam}.

Similarly, we obtain a pairing on de Rham cohomology which is normalized by $(2\pi i)^{-r}$. For a class
$x$ at weight $\mathbf a$ and a class $y$ at weight $-\mathbf a$,
we use the argument order of \cite[equation (15.1)]{Lam} and write
\[
 \langle x,y\rangle^{\nabla}_{\mathbf a}
 =\frac{1}{(2\pi i)^r}\int_U y\wedge\reg(x).
\]
Here the integral uses a compactly supported smooth representative
of $\reg(x)$.

Lam's comparison theorems~\cite[Theorems 12.9 and 12.12]{Lam} identify the geometric pairings with Lam's combinatorial
Betti and de Rham pairings:
\begin{align*}
 \langle C_{-\mathbf a}(P),C_{\mathbf a}(Q)\rangle_{\mathrm{hom}}
   &=B_b^{\cB}(P,Q),\\
 \langle\iota_*x,\iota_*y\rangle^{\nabla}_{\mathbf a}
   &=\langle x,y\rangle^{\dR}_{\mathbf a}.
\end{align*}
Here $P,Q\in\cB$, $x\in A_{\mathbf a}$ and $y\in A_{-\mathbf a}$.
 In particular,
$D_{\mathbf a}^{\mathrm{nbc}}$ and $D_{\mathbf a}^{\mathrm{ch}}$ are also the Gram
matrices of the normalized geometric de Rham pairing in the two
displayed bases.

Under the perfect evaluation pairings $(~,~)$, let
$\{\xi_{\pm\mathbf a}(P)\}_{P\in\cB}$ be the bases of
$H^r(U,\Loc_{\pm\mathbf a})$ dual to
$\{\reg C_{\mp\mathbf a}(P)\}_{P\in\cB}$, respectively.
Passing to dual bases takes the inverse transpose Gram
matrix \cite[Section 12.4]{Lam}, so symmetry of $B_b^{\cB}$ gives
\begin{equation}\label{eq:dual-Gram}
 \langle\xi_{\mathbf a}(P),\xi_{-\mathbf a}(Q)\rangle_{\mathrm{coh}}
 =\bigl((B_b^{\cB})^{-1}\bigr)(P,Q).
\end{equation}

\subsection{The period relation}\label{subsec:period}
For a real arrangement and nonresonant weights, define the regularized
$\cB\times\Bnbc$ period matrix by
\[
 \PM^*_{\mathbf a}(P,B)
 =\bigl(\per_{\mathbf a}(\psi_{\mathbf a}^B),\reg C_{-\mathbf a}(P)\bigr).
\]
Thus rows of the period matrix are indexed by bounded chambers and columns by $\beta$nbc bases of $M$.
The parentheses denote the evaluation pairing between cohomology
and homology with dual local coefficients. This is the notation of
\cite[Definition 3.3.2]{DouaiTerao1997}.
Since the $\xi_{\mathbf a}(P)$ are evaluation-dual to the regularized
chamber cycles, we have
\begin{equation}\label{eq:period-coordinates}
 \per_{\mathbf a}(\psi_{\mathbf a}^B)
 =\sum_{P\in\cB}\PM^*_{\mathbf a}(P,B)\,\xi_{\mathbf a}(P).
\end{equation}


For $\operatorname{Re}a_i>0$, $i\in E_0$, the chamber integrals
converge and define the matrix
\[
 \PM_{\mathbf a}(P,B)
 =\int_P\varphi_{\mathbf a,P}\,\iota(\psi_{\mathbf a}^B).
\]
At nonresonant weights in this region, $\PM_{\mathbf a}=\PM^*_{\mathbf a}$
\cite[Remark 3.3.3]{DouaiTerao1997}.
The matrix $\PM^*_{-\mathbf a}$ uses the opposite loading
and the forms $\psi_{-\mathbf a}^B$.

\begin{example}\label{ex:line3}
Continuing Examples~\ref{ex:line} and~\ref{ex:line2}, the form
$\psi^B_{\mathbf a}=a_2\bar e_2$ is realized as
$\iota(\psi^B_{\mathbf a})=-a_2\,\dd x/(1-x)$.
For $\operatorname{Re}a_1,\operatorname{Re}a_2>0$, its period is
\[
 \PM_{\mathbf a}(P,B)
 =-a_2\int_0^1 x^{a_1}(1-x)^{a_2-1}\,\dd x
 =-a_2 B(1+a_1,a_2),
\]
where $B(x,y)$ is the classical beta function. The expression
\[
 \PM_{\mathbf a}^*(P,B)
 =-\frac{\Gamma(1+a_1)\Gamma(1+a_2)}{\Gamma(1-a_0)}
\]
gives its meromorphic continuation.
\end{example}


\begin{proposition}\label{prop:period-relation}
For nonresonant weights with $b_i=e^{-\pi i a_i}$,
\begin{equation}\label{eq:period-relation}
 (2\pi i)^rD_{\mathbf a}^{\mathrm{nbc}}
 =(\PM^*_{\mathbf a})^{\mathsf T}(B_b^{\cB})^{-1}\PM^*_{-\mathbf a}.
\end{equation}
Consequently,
\begin{equation}\label{eq:det-period-relation}
 \det B_b^{\cB}=(2\pi i)^{-r\beta(M)}
       \frac{\det\PM^*_{\mathbf a}\det\PM^*_{-\mathbf a}}
            {\det D_{\mathbf a}^{\mathrm{nbc}}}.
\end{equation}
\end{proposition}
\begin{proof}
The comparison theorem gives
\[
 (2\pi i)^r\langle x,y\rangle^{\dR}_{\mathbf a}
 =\langle\per_{\mathbf a}(x),\per_{-\mathbf a}(y)\rangle_{\mathrm{coh}}
\]
by \cite[equation (15.1)]{Lam}. Substituting
\eqref{eq:period-coordinates} at both weights and using
\eqref{eq:dual-Gram} gives \eqref{eq:period-relation}.
This is the twisted period relation of Cho--Matsumoto
\cite{ChoMatsumoto1995} in the form of
\cite[equation (15.2)]{Lam}. Taking determinants gives
\eqref{eq:det-period-relation}.
\end{proof}

\subsection{Period determinants and continuation}\label{subsec:period-determinant}
For a bounded region $P$ and $i\in E_0$, let $\Sigma_{i,P}$ be
the face of $\bar P$ on which $|f_i|$ is maximal. Set
\[
 c_{i,P}=|f_i(\Sigma_{i,P})|>0,\qquad
 C(\mathbf a)=\prod_{i\in E_0}\prod_{P\in\cB}c_{i,P}^{a_i}.
\]
These are the external supporting-face critical values in
\cite[Sections 2.6 and 4.3.2]{MukherjeeVarchenko2009}, expressed in
our positive normalization. In particular,
\begin{equation}\label{eq:critical-cancellation}
 C(\mathbf a)C(-\mathbf a)=1.
\end{equation}
The beta function of the weighted arrangement is
\[
 \Betafunc_M(\mathbf a)
 =\frac{\displaystyle
      \prod_{F\in\cL\setminus(\cL_0\cup\{\varnothing\})}
          \Gamma(1+a_F)^{w_F}}
        {\displaystyle
      \prod_{F\in\cL_0\setminus\{E\}}\Gamma(1-a_F)^{w_F}}.
\]

\begin{theorem}[Loeser--Sabbah, Douai--Terao]\label{thm:period-determinant}
For $\operatorname{Re}a_i>0$, $i\in E_0$,
\[
 \det\PM_{\mathbf a}=\varepsilon\, C(\mathbf a)\Betafunc_M(\mathbf a),
 \qquad\varepsilon\in\{1,-1\},
\]
where $\varepsilon$ depends only on the fixed orderings and orientations.
\end{theorem}
\begin{proof}
We use \cite{LoeserSabbah1991} and
\cite[Theorem 5.1.1]{DouaiTerao1997} in the formulation of
\cite[Theorem 4.1 and Sections 2.5, 4.2.1, 4.3]{MukherjeeVarchenko2009}.
The forms \eqref{eq:nbc-form} are their weighted flag forms, the
exponents are $w_F$, and the critical-value product for the loading
\eqref{eq:loading} is $C(\mathbf a)$. Changing the orderings or
orientations contributes the fixed sign $\varepsilon$.
\end{proof}

The entries of $\PM_{\mathbf a}$ are holomorphic for
$\operatorname{Re}a_i>0$, $i\in E_0$. Their meromorphic continuations
to the full weight space $a_E=0$ are the entries of $\PM^*_{\mathbf a}$
\cite[Remark 3.3.3]{DouaiTerao1997}.
Since the determinant is a polynomial in the entries,
Theorem~\ref{thm:period-determinant} therefore gives
\[
 \det\PM^*_{\mathbf a}=\varepsilon\,C(\mathbf a)\Betafunc_M(\mathbf a)
\]
as an identity of meromorphic functions on $a_E=0$. We can now
apply it at both $\mathbf a$ and $-\mathbf a$, using the geometric
comparisons at nonresonant weights.

\begin{corollary}\label{cor:opposite-periods}
As meromorphic functions,
\[
 \det\PM^*_{\mathbf a}\det\PM^*_{-\mathbf a}
 =\Betafunc_M(\mathbf a)\Betafunc_M(-\mathbf a).
\]
\end{corollary}
\begin{proof}
Multiply the continued formulas at $\mathbf a$ and $-\mathbf a$.
The critical-value factors cancel by \eqref{eq:critical-cancellation},
and the signs multiply to $\varepsilon^2=1$.
\end{proof}

\begin{remark}
Using $\varphi_{\mathbf a}$ in place of the normalized branches in
\eqref{eq:loading} replaces $B_b^{\cB}$ by
$\Lambda B_b^{\cB}\Lambda^{-1}$, where
$\Lambda=\diag(b_{P^-})_{P\in\cB}$.
The period matrices $\PM^*_{\mathbf a}$ and $\PM^*_{-\mathbf a}$
are multiplied on the left by $\Lambda$ and $\Lambda^{-1}$,
respectively. Thus the Betti determinant and the product of period
determinants are unchanged.
\end{remark}
